\chapter{Architektura aplikacji}
\label{chap:architektura}

\section{Ogólna architektura systemu}

Aplikacja \textit{IdentityCore} została zaprojektowana jako lokalna aplikacja desktopowa, w której wydzielono kilka głównych obszarów odpowiedzialności: warstwę prezentacji, logikę aplikacyjną, moduły biometryczne oraz warstwę dostępu do danych. Taki podział pozwala oddzielić elementy odpowiedzialne za wygląd aplikacji od części realizujących operacje identyfikacyjne i zapis danych.

Na poziomie architektury system można traktować jako aplikację zorganizowaną warstwowo. Widoki odpowiadają za prezentację danych, modele widoków za obsługę stanu ekranów i akcji operatora, serwisy za logikę aplikacyjną, a katalog \texttt{Data} za komunikację z bazą danych. Moduły biometryczne są częścią logiki aplikacyjnej, ale zostały wydzielone tak, aby można było rozwijać je niezależnie od warstwy interfejsu.

\section{Podział na warstwy i moduły}

Najważniejsze warstwy oraz moduły aplikacji przedstawiono w tabeli~\ref{tab:warstwy-moduly}. Zestawienie pokazuje, za co odpowiadają poszczególne części projektu i gdzie znajdują się ich przykładowe elementy.

\begin{xltabular}{\textwidth}{|p{2.9cm}|X|p{5.8cm}|}
    \hline
    \textbf{Moduł/warstwa} & \textbf{Odpowiedzialność} & \textbf{Przykładowe elementy \newline projektu} \\
    \hline
    \endfirsthead

    \hline
    \textbf{Moduł/warstwa} & \textbf{Odpowiedzialność} & \textbf{Przykładowe elementy \newline projektu} \\
    \hline
    \endhead

    Warstwa widoków & Prezentacja interfejsu użytkownika oraz obsługa ekranów aplikacji. & Katalog \texttt{Views}, między innymi widoki panelu głównego, rejestru osób, \newline identyfikacji, historii, ustawień i \newline logowania. \\
    \hline
    Modele widoków & Przechowywanie stanu ekranów, obsługa komend oraz \newline pośredniczenie pomiędzy \newline widokiem a logiką aplikacji. & Katalog \texttt{ViewModels}, między innymi modele widoków panelu głównego, \newline rejestru osób, identyfikacji, historii i \newline ustawień. \\
    \hline
    Serwisy aplikacyjne & Realizacja operacji wymagających osobnej logiki aplikacyjnej lub przetwarzania danych. & Katalog \texttt{Services}, między innymi \newline serwisy logowania, kamery, \newline przetwarzania obrazów, identyfikacji i ustawień. \\
    \hline
    Warstwa dostępu \newline do danych & Odczyt, zapis i aktualizacja danych przechowywanych przez aplikację. & Katalog \texttt{Data}, między innymi \newline repozytoria osób, szablonów \newline biometrycznych, historii prób, panelu głównego oraz połączenia z bazą. \\
    \hline
    Modele danych & Reprezentacja obiektów \newline wykorzystywanych w systemie oraz wyników operacji. & Katalog \texttt{Models}, między innymi modele osoby, próbek biometrycznych, wyników identyfikacji, kandydatów dopasowania i ustawień decyzyjnych. \\
    \hline
    Moduł identyfikacji \newline twarzy & Obsługa procesu związanego z \newline analizą i porównywaniem danych twarzy. & Elementy w katalogach \texttt{Services}, \texttt{Models}, \texttt{Data}, \texttt{Assets} oraz \texttt{MLModels}. \\
    \hline
    Moduł identyfikacji \newline odcisku palca & Obsługa procesu związanego z \newline analizą i porównywaniem danych odcisku palca. & Elementy w katalogach \texttt{Services}, \texttt{Models} oraz \texttt{Data}. \\
    \hline
    Moduł historii prób & Przechowywanie i udostępnianie informacji o wykonanych próbach identyfikacji. & Widok historii, model widoku historii, repozytorium identyfikacji oraz model próby identyfikacji. \\
    \hline
    Moduł ustawień \newline systemowych & Obsługa parametrów \newline wpływających na działanie procesu identyfikacji. & Widok ustawień, model widoku \newline ustawień, \texttt{SystemSettingsService} oraz model ustawień biometrycznych. \\
    \hline

    \caption{Główne warstwy i moduły aplikacji}
    \label{tab:warstwy-moduly}\\
\end{xltabular}

Podział przedstawiony w tabeli porządkuje odpowiedzialności w projekcie. Najważniejsze jest to, że widoki nie wykonują samodzielnie operacji biometrycznych ani nie obsługują bezpośrednio zapisu danych. Zadania te są przekazywane do odpowiednich modeli widoków, serwisów lub klas z warstwy \texttt{Data}.

\newpage

\section{Struktura projektu i organizacja kodu}

Struktura katalogów pokazuje fizyczne rozmieszczenie poszczególnych elementów aplikacji w projekcie. Uproszczony układ katalogów przedstawiono poniżej.

\begin{verbatim}
IdentityCore.App/
|-- Assets/
|   `-- Cascades/
|-- Data/
|-- Database/
|-- MLModels/
|-- Models/
|-- Services/
|-- Styles/
|-- ViewModels/
|-- Views/
|-- App.xaml
|-- App.xaml.cs
|-- AssemblyInfo.cs
|-- IdentityCore.App.csproj
|-- MainWindow.xaml
`-- MainWindow.xaml.cs
\end{verbatim}

Katalog \texttt{Assets} przechowuje katalog \texttt{Cascades} zawierający pliki kaskad wykorzystywane pomocniczo podczas pracy z obrazami twarzy.

Katalog \texttt{Data} zawiera klasy odpowiedzialne za dostęp do danych. W tym miejscu umieszczono połączenie z bazą danych oraz repozytoria obsługujące dane osób, próbek biometrycznych, szablonów twarzy, szablonów odcisków palców, historii identyfikacji, panelu głównego i logowania.

Katalog \texttt{Database} zawiera skrypty SQL służące do tworzenia, aktualizacji oraz czyszczenia danych biometrycznych w bazie. Szczegółowy opis struktury danych zostanie przedstawiony w osobnym rozdziale dotyczącym bazy danych.

Katalog \texttt{MLModels} przechowuje plik modelu wykorzystywany przez moduł identyfikacji twarzy. Oddzielenie tego typu zasobów od kodu aplikacji ułatwia utrzymanie porządku w projekcie.

Katalog \texttt{Models} zawiera klasy opisujące dane wykorzystywane w systemie, między innymi osoby, próbki biometryczne, wyniki identyfikacji, kandydatów dopasowania, szablony oraz ustawienia decyzyjne.

Katalog \texttt{Services} grupuje logikę aplikacyjną, która nie należy bezpośrednio do widoków. Obejmuje to między innymi obsługę logowania, kamery, przetwarzania obrazów, identyfikacji, diagnostyki, kalibracji i ustawień.

Katalog \texttt{Styles} zawiera współdzielone pliki XAML odpowiedzialne za wygląd elementów interfejsu, takich jak przyciski, karty, kolory, tabele i pola wejściowe.

Katalogi \texttt{ViewModels} i \texttt{Views} odpowiadają za warstwę użytkową aplikacji. Pierwszy z nich przechowuje modele widoków obsługujące stan ekranów i komendy, a drugi zawiera pliki widoków oraz powiązane z nimi pliki code-behind. Pliki \texttt{App.xaml} i \texttt{MainWindow.xaml} odpowiadają za uruchomienie aplikacji oraz główne okno programu.

\section{Zastosowane wzorce i zasady projektowe}

Organizacja aplikacji opiera się na podejściu MVVM. W projekcie oznacza to, że widok odpowiada za prezentację danych, a model widoku obsługuje stan ekranu, komendy i przekazywanie zadań do dalszych warstw aplikacji. Dzięki temu logika działania nie musi być umieszczana bezpośrednio w plikach widoków.

Serwisy odpowiadają za operacje wymagające osobnej logiki aplikacyjnej. Dotyczy to między innymi przetwarzania obrazów, obsługi kamery, logowania, ustawień oraz identyfikacji biometrycznej. Wydzielenie serwisów ogranicza rozrastanie się modeli widoków i pozwala utrzymać bardziej czytelny podział zadań.

Najważniejszą zasadą projektową jest separacja odpowiedzialności. Każdy element aplikacji ma określoną rolę: widok prezentuje dane, model widoku obsługuje zachowanie ekranu, serwis wykonuje logikę aplikacyjną, model opisuje dane, a warstwa \texttt{Data} odpowiada za ich trwały zapis i odczyt.

\section{Walidacja danych}

Walidacja danych dotyczy sprawdzenia poprawności informacji przed wykonaniem operacji. Jej celem jest ograniczenie sytuacji, w których aplikacja próbuje zapisać niepełne dane osoby, uruchomić identyfikację bez próbki wejściowej albo wykonać porównanie bez danych odniesienia.

W rejestrze osób walidacja obejmuje kontrolę wymaganych pól i podstawowej poprawności danych profilu. W modułach biometrycznych sprawdzane jest przede wszystkim to, czy dostępny jest plik lub obraz możliwy do dalszej analizy. Jeżeli dane są niekompletne, operacja nie powinna być wykonywana do momentu ich poprawienia.

Walidacja pełni więc funkcję zabezpieczenia przed typowymi błędami użytkowania aplikacji. Szczegółowy sposób prezentowania komunikatów operatorowi zostanie opisany w rozdziale dotyczącym warstwy użytkowej.

\section{Obsługa błędów i wyjątków}

Obsługa błędów dotyczy sytuacji, w których mimo sprawdzenia danych operacja nie może zostać wykonana poprawnie. Mogą to być problemy z połączeniem z bazą danych, niepowodzenie zapisu lub odczytu, niepoprawny plik wejściowy, błąd przetwarzania obrazu albo nieudana próba logowania.

W takich przypadkach aplikacja powinna reagować w sposób kontrolowany. Typowy problem nie powinien powodować nagłego przerwania działania programu, lecz zakończyć się przekazaniem operatorowi informacji o przyczynie niepowodzenia. Pozwala to poprawić dane, wybrać inny plik, ponowić próbę albo sprawdzić konfigurację środowiska.

Obsługa wyjątków uzupełnia walidację. Walidacja zapobiega części błędów przed rozpoczęciem operacji, natomiast obsługa wyjątków odpowiada za przypadki, których nie dało się przewidzieć lub które wystąpiły w trakcie działania aplikacji.

\section{Czytelność i utrzymanie projektu}

Przyjęty podział projektu ułatwia analizę i dalszą rozbudowę aplikacji. Najważniejsze znaczenie ma oddzielenie widoków od logiki aplikacyjnej, wydzielenie serwisów oraz skupienie dostępu do danych w katalogu \texttt{Data}. Dzięki temu projekt pozostaje czytelny i możliwy do rozwijania w ramach kolejnych etapów pracy.